\clearpage\fancyhead[L]{\small JUNO.}\hypertarget{native-375-lower}{}\label{native-375-lower}
\begin{center}\textbf{JUNO.}\par\vspace{4pt}\hrule\vspace{4pt}\fontsize{10}{12}\selectfont Monatliche Correspondenz. Vol. XXVIII. pp. 574\;..\;578. December 1813.\par\vspace{4pt}\hrule\vspace{4pt}\end{center}


\begin{flushright}Göttingen, 8, 9 and 16 December 1813.\end{flushright}
\noindent --- --- --- The weather has been very unfavorable for observing, as usual at this time of year. I now determine the time from stellar azimuths, which I observe (in daylight) with the theodolite. Observations at equal solar altitudes never succeed at present, \SourceLeaf{376}since the sky clears a little only for isolated hours now and then. Toward sunset I now compare $\alpha$ Aurigae with a tower 4000 meters away whose azimuth was determined in the summer, and that works quite well.

\Needspace{12\baselineskip}\par\indent I was able to observe Juno only once, but the observation itself is very good and was made on the very day of the opposition. The planet was of the eighth magnitude.
\begin{center}\fontsize{10}{12}\selectfont\setlength{\tabcolsep}{4pt}\renewcommand{\arraystretch}{1.15}\begin{tabular}{l|c|c|c}
1813 & \shortstack{Mean time\\in Göttingen} & \shortstack{Right\\ascension $\JunoSymbol$} & \shortstack{Southern\\declination $\JunoSymbol$}\\\hline
Nov. 19 & $9^{\mathrm u}\;46^{\mathrm m}\;46^{\mathrm s}$ & $60^\circ\;30'\;35''0$ & $3^\circ\;5'\;9''3$\\
\end{tabular}\end{center}
\Needspace{7\baselineskip}\par\indent The apparent position of the comparison star was determined from the \emph{Hist. Cél.}
\begin{center}R.A. $=60^\circ\;47'\;50''0$\quad Decl. $=3^\circ\;0'\;16''0$ S.\end{center}

\Needspace{12\baselineskip}\par\indent I derived the opposition from this observation without any hesitation, since observations are hardly likely to come in from elsewhere.
\begin{center}\emph{Eighth opposition of Juno:}\end{center}
\begin{center}1813 November 19. $18^{\mathrm u}\;11^{\mathrm m}\;25^{\mathrm s}$ Mean time in Göttingen\end{center}
\begin{center}\fontsize{10}{12}\selectfont\setlength{\tabcolsep}{2pt}\renewcommand{\arraystretch}{1.15}\begin{tabular}{lcr}
True longitude & $.\;.\;.\;.\;.\;.\;.\;.\;.\;.\;.\;.\;.$ & $57^\circ\;33'\;58''0$\\
Geocentric latitude & $.\;.\;.\;.\;.\;.\;.\;.\;.\;.$ & $23\;18\;46.8\ \text{S.}$\\
\end{tabular}\end{center}
\Needspace{6\baselineskip}\begin{center}\emph{Continuation of 9 December $10^{\mathrm u}$.}\end{center}
\par\indent This evening, in thick fog and moonlight, I still did not check again the place where I observed Juno on 19 November, but I did seek out Juno itself again. It seemed \emph{relatively} \emph{brighter} today than on 19 November. At least it was considerably brighter than the eighth-magnitude star it was compared with today, whereas on 19 November it scarcely matched the eighth-magnitude star used for comparison. I ask Mr. \textsc{Nicolai} to reduce the observations and compare them with the elements given below. The observation is based on only three poor comparisons with a star from the \emph{Hist. Cél.}, which was observed at the second wire on 20 January 1798 at $3^{\mathrm u}\;42^{\mathrm m}\;48^{\mathrm s}5$. Even so, it may serve as a check.
\begin{center}
1813 December 9. $0^{\mathrm u}\;22^{\mathrm m}\;2^{\mathrm s}$ sidereal time, $\JunoSymbol$ follows the star
$3^{\mathrm m}\;49^{\mathrm s}75$ of time $=57'\;26''25$ of arc

$0^{\mathrm u}\;32^{\mathrm m}\;4^{\mathrm s}$ --- $\JunoSymbol$ is farther south by $5'\;0''9$.
\end{center}
\noindent This gives approximately
\begin{center}R.A. $=56^\circ\;56\tfrac12'$\quad Decl. $=3^\circ\;56\tfrac12'$ S.\end{center}

\Needspace{15\baselineskip}\par\indent The elements based on the oppositions of 1810, 1811, 1812, 1813 are as follows (X):
\begin{center}\fontsize{9}{11}\selectfont\setlength{\tabcolsep}{2pt}\renewcommand{\arraystretch}{1.15}\begin{tabular}{lcr}
Epoch 1814, Göttingen meridian & $.\;.\;.\;.\;.\;.$ & $65^\circ\;19'\;23''8$\\
Daily tropical motion & $.\;.\;.\;.\;.\;.\;.\;.\;.$ & $812''709$\\
Log. of the semimajor axis & $.\;.\;.\;.\;.\;.\;.\;.\;.\;.$ & $0.4267631$\\
Perihelion (1814) & $.\;.\;.\;.\;.\;.\;.\;.\;.\;.\;.\;.\;.\;.$ & $53^\circ\;10'\;0''40$\\
Angle of eccentricity & $.\;.\;.\;.\;.\;.\;.\;.\;.\;.\;.\;.\;.$ & $14\;43\;7.97$\\
Ascending node (1814) & $.\;.\;.\;.\;.\;.\;.\;.\;.\;.$ & $171\;10\;33.00$\\
Inclination & $.\;.\;.\;.\;.\;.\;.\;.\;.\;.\;.\;.\;.\;.\;.\;.\;.$ & $13\;4\;12.03$\\
\end{tabular}\end{center}
\par\indent Comparison of these elements, especially the mean motion, with the elements (VIII) based on the oppositions of 1804, 1806, 1807, 1808 shows how strongly Jupiter’s perturbations make themselves felt.

\Needspace{18\baselineskip}\begin{flushright}\leavevmode\SourceLeaf{377}Göttingen, 16 December 1813.\end{flushright}
\noindent --- --- --- I did take the trouble after all to fit the elements to the latest opposition of Pallas, IX, as well. The corrections are extremely small:
\begin{center}\fontsize{10}{12}\selectfont\setlength{\tabcolsep}{2pt}\renewcommand{\arraystretch}{1.15}\begin{tabular}{lcr}
Longitude of the node & $.\;.\;.\;.\;.\;.$ & $-2''7$\\
Inclination of the orbit & $.\;.\;.\;.\;.\;.$ & $-1.4$\\
Longitude of perihelion & $.\;.\;.\;.\;.\;.$ & $-9.4$\\
Angle of eccentricity & $.\;.\;.\;.\;.\;.$ & $+1.8$\\
Epoch 1810 & $.\;.\;.\;.\;.\;.\;.$ & $-4.2$\\
Daily motion & $.\;.\;.\;.$ & $-0''00569$\\
Log. of the semimajor axis & $.$ & $+0.0000021$\\
\end{tabular}\end{center}
\Needspace{22\baselineskip}\par\indent The agreement of the 9 oppositions with the refined elements is now as follows:
\begin{center}Difference:\end{center}
\begin{center}\fontsize{10}{12}\selectfont\setlength{\tabcolsep}{7pt}\renewcommand{\arraystretch}{1.15}\begin{tabular}{l|c|c}
Oppos. & \shortstack{Mean\\longitude} & \shortstack{Heliocentric\\latitude.}\\\hline
I & $-1''3$ & $-5''3$\\
II & $-12.9$ & $-11.9$\\
III & $-20.4$ & $-0.4$\\
IV & $+2.5$ & $-6.8$\\
V & $+23.5$ & $+9.0$\\
VI & $+21.5$ & $-3.1$\\
VII & $+14.1$ & $-7.2$\\
VIII & $-16.1$ & $+4.0$\\
IX & $-10.9$ & $-0.0$\\
X & $-14.0$ & $+2.4$\\
\end{tabular}\end{center}
\par\indent I am very curious to see how the differences will turn out once Saturn’s perturbations are included. [The entry for opposition X is a handwritten remark. \textsc{Schering}.]
